\chapter{จุลินทรีย์กับการเปลี่ยนแปลงสภาพภูมิอากาศและการเกษตรยั่งยืน}
\label{chap:ch08}

% ==========================================================
% แผนบริหารการสอนประจำบทที่ 8
% ==========================================================
\section*{แผนบริหารการสอนประจำบทที่ 8}
\addcontentsline{toc}{section}{แผนบริหารการสอนประจำบทที่ 8}

\noindent\textbf{หัวข้อเนื้อหาประจำบท}
\begin{enumerate}
    \item ภาคการเกษตรกับวิกฤตการณ์การปล่อยก๊าซเรือนกระจก ($CH_4, N_2O, CO_2$)
    \item กลไกของจุลินทรีย์ในการสร้างและการลดการปล่อยก๊าซเรือนกระจก
    \item การกักเก็บคาร์บอนในดิน (Soil Carbon Sequestration) และชีวมวลซากจุลินทรีย์ (Necromass)
    \item ปฏิสัมพันธ์เชิงส่งเสริมระหว่างถ่านชีวภาพ (Biochar) กับจุลินทรีย์ดิน
    \item การพัฒนาจุลินทรีย์ทนแล้งทนเค็มเพื่อรับมือกับภาวะสภาพอากาศสุดขั้ว
\end{enumerate}

\noindent\textbf{วัตถุประสงค์เชิงพฤติกรรม}
\begin{enumerate}
    \item วิเคราะห์บทบาทของจุลินทรีย์ต่อการปลดปล่อยและดูดซับก๊าซเรือนกระจกในแปลงเกษตรได้
    \item อธิบายกลไกการเปลี่ยนสารอินทรีย์ให้กลายเป็นคาร์บอนที่เสถียรในดินผ่าน Necromass ได้
    \item อธิบายคุณสมบัติของ Biochar ในการเป็นแหล่งพักพิงและเพิ่มความอยู่รอดของจุลินทรีย์ดินได้
    \item ออกแบบแนวทางการจัดการดินเพื่อลดการปล่อยก๊าซมีเทนและไนตรัสออกไซด์ในพื้นที่เกษตรได้
\end{enumerate}

\noindent\textbf{กิจกรรมการเรียนการสอน}
\begin{enumerate}
    \item การบรรยายและวิเคราะห์ข้อมูลฟลักซ์ก๊าซเรือนกระจกในระบบนิเวศเกษตร
    \item การทดลองจำลองการตรึงจุลินทรีย์ลงในถ่านชีวภาพและประเมินอัตราการรอดชีวิต
    \item การอภิปรายแนวทางการทำนาข้าวแบบเปียกสลับแห้ง (AWD) เพื่อลดก๊าซมีเทน
\end{enumerate}

\noindent\textbf{สื่อการเรียนการสอน}
\begin{enumerate}
    \item เอกสารคำสอนบทที่ 8 และแผนภาพแสดงสมดุลคาร์บอนในดิน
    \item ตัวอย่างถ่านชีวภาพจากแกลบข้าวและกะลาทุเรียน
\end{enumerate}

\noindent\textbf{การประเมินผล}
\begin{enumerate}
    \item การประเมินผลจากการเขียนข้อเสนอโครงการจัดการดินเพื่อคาร์บอนเครดิตภาคเกษตร
    \item การทดสอบความรู้เชิงลึกเกี่ยวกับวัฏจักรคาร์บอนและก๊าซเรือนกระจก
\end{enumerate}

\newpage

% ==========================================================
% เนื้อหาบทเรียน
% ==========================================================

\section{ภาคการเกษตรกับวิกฤตสภาพภูมิอากาศ}
\index{การเปลี่ยนแปลงสภาพภูมิอากาศ (Climate Change)}

ภาคการเกษตรเป็นทั้งแหล่งปล่อยก๊าซเรือนกระจก (Greenhouse Gas Emissions) ที่สำคัญของโลก และในขณะเดียวกันก็เป็นภาคส่วนที่ได้รับผลกระทบจากความแปรปรวนของสภาพอากาศอย่างรุนแรงที่สุด ก๊าซเรือนกระจกหลัก 2 ชนิดที่เกิดจากกิจกรรมของจุลินทรีย์ในดินเกษตร ดังนี้
\begin{enumerate}
    \item \textbf{ก๊าซมีเทน ($CH_4$)} เกิดขึ้นจากการย่อยสลายสารอินทรีย์ของกลุ่มอาร์เคียเมทาโนเจน (Methanogens) ภายใต้สภาวะขาดออกซิเจนในแปลงนาข้าวที่มีน้ำท่วมขังขังยาวนาน และในระบบทางเดินอาหารของสัตว์เคี้ยวเอื้อง
    \item \textbf{ก๊าซไนตรัสออกไซด์ ($N_2O$)} เกิดขึ้นจากกระบวนการไนตริฟิเคชันและดีไนตริฟิเคชันที่ไม่สมบูรณ์ เมื่อมีการใส่ปุ๋ยไนโตรเจนเคมีเกินความต้องการของพืช
\end{enumerate}

\section{การกักเก็บคาร์บอนในดินและชีวมวลซากจุลินทรีย์}
\index{การกักเก็บคาร์บอนในดิน (Soil Carbon Sequestration)}

มุมมองดั้งเดิมทางปฐพีศาสตร์เชื่อว่า สารฮิวมัสที่เสถียรในดินเกิดจากการตกค้างของสารประกอบโครงสร้างแข็ง เช่น ลิกนินจากพืช ทว่างานวิจัยทางชีววิทยาโมเลกุลในยุคปัจจุบันพิสูจน์แล้วว่า มากกว่า 50--80\% ของอินทรีย์คาร์บอนที่เสถียรในดิน (Stable Soil Organic Carbon) มีที่มาจาก \textbf{ชีวมวลซากจุลินทรีย์ที่ตายแล้ว (Microbial Necromass)}

เมื่อเซลล์ของแบคทีเรียและเชื้อราตายลง ผนังเซลล์ที่ประกอบด้วยเปปทิโดไกลแคน กลูแคน และไคติน จะแตกตัวออกและเข้าทำอันตรกิริยาทางเคมีกับอนุภาคดินเหนียวและออกไซด์ของเหล็ก/อะลูมิเนียม เกิดเป็นสารประกอบเชิงซ้อนระหว่างแร่ธาตุกับสารอินทรีย์ (Mineral-Associated Organic Matter: MAOM) ซึ่งมีความคงทนต่อการย่อยสลายสูงมาก สามารถกักเก็บคาร์บอนไว้ในดินได้ยาวนานนับศตวรรษ

\begin{figure}[H]
\centering
\begin{tikzpicture}[scale=0.95]
    % Plant aboveground
    \draw[line width=3pt, color=agriLeaf] (5, 3.5) -- (5, 0.5);
    \draw[fill=agriLeaf!80!white, draw=agriGreen] (4.2, 2.5) .. controls (3.5, 3.5) .. (5, 3.8) .. controls (6.5, 3.5) .. (5.8, 2.5) -- cycle;
    \node at (5, 4.2) [font=\footnotesize\bfseries\color{agriGreen}] {การสังเคราะห์ด้วยแสง ตรึง $CO_2$};

    % Soil surface
    \draw[line width=2pt, color=agriEarth] (0, 0.5) -- (10.5, 0.5);
    \node at (1.0, 0.8) [font=\tiny\bfseries\color{agriEarth}] {ผิวดิน};

    % Rhizodeposition arrow down
    \draw[-{Stealth[scale=1.3]}, line width=2pt, color=agriGold] (5, 0.5) -- (5, -1.0)
        node[midway, right, font=\tiny\bfseries\color{agriGold!80!black}] {สารหลั่งรากและซากพืช};

    % Living Microbial Biomass
    \node (microbes) [rectangle, rounded corners=6pt, draw=agriGreen, fill=agriMint, line width=1.5pt, minimum width=4.5cm, minimum height=1.0cm, align=center, font=\footnotesize\bfseries, below=1.0cm of {5,0}]
        {ชีวมวลจุลินทรีย์ที่มีชีวิต (Living Biomass)\\แบคทีเรียและเชื้อราเปลี่ยน $C$ เป็นเซลล์};

    % Arrows to respiration and necromass
    \draw[-{Stealth[scale=1.2]}, line width=1.4pt, color=agriRed, dashed] (microbes.west) .. controls (2.0, -1.5) .. (2.0, 1.2)
        node[midway, left, align=center, font=\tiny\bfseries\color{agriRed}] {หายใจปลดปล่อย\\$CO_2$ ($40-60\%$)};

    \node (necromass) [rectangle, rounded corners=6pt, draw=agriEarth, fill=agriEarth!20, line width=1.5pt, minimum width=4.5cm, minimum height=1.0cm, align=center, font=\footnotesize\bfseries, below=1.5cm of microbes]
        {ซากเซลล์จุลินทรีย์ (Microbial Necromass)\\ผนังเซลล์ Peptidoglycan \& Chitin};

    \draw[-{Stealth[scale=1.4]}, line width=2pt, color=agriGreen] (microbes) -- (necromass)
        node[midway, right, font=\tiny\bfseries\color{agriGreen}] {การหมุนเวียนเซลล์ (Cell Turnover)};

    % Stable MAOM pool
    \node (maom) [rectangle, rounded corners=8pt, draw=agriBlue, fill=agriBlue!15, line width=2pt, minimum width=6.5cm, minimum height=1.2cm, align=center, font=\small\bfseries\color{agriBlue!80!black}, below=1.2cm of necromass]
        {คาร์บอนเสถียรยึดเกาะแร่ดิน (MAOM)\\กักเก็บคาร์บอนยาวนาน $> 100$ ปี};

    \draw[-{Stealth[scale=1.4]}, line width=2pt, color=agriBlue] (necromass) -- (maom)
        node[midway, right, font=\tiny\bfseries\color{agriBlue}] {จับกับแร่ดินเหนียว (Mineral Binding)};
\end{tikzpicture}
\caption{วิถีการกักเก็บคาร์บอนในดินผ่านชีวมวลซากจุลินทรีย์ (Microbial Carbon Pump and Necromass)}
\label{fig:carbon_pump}
\end{figure}

\section{การเสริมฤทธิ์ระหว่างถ่านชีวภาพกับจุลินทรีย์}
\index{ถ่านชีวภาพ (Biochar-Microbe Interactions)}

\textbf{ถ่านชีวภาพ (Biochar)} ผลิตจากการเผาวัสดุเหลือทิ้งทางการเกษตร (แกลบ กะลามะพร้าว กากตะกอนอ้อย) ภายใต้สภาวะจำกัดออกซิเจน (Pyrolysis) มีโครงสร้างรูพรุนระดับนาโนและไมโครสูงมาก มีพื้นที่ผิวจำเพาะมากกว่า 200--400 ตารางเมตรต่อกรัม การใส่ถ่านชีวภาพร่วมกับหัวเชื้อจุลินทรีย์ส่งผลดีอย่างยิ่ง ดังนี้
\begin{itemize}
    \item เป็นแหล่งพักพิงและปกป้องจุลินทรีย์ (Microhabitat refuge) จากการถูกล่าโดยโพรโทซัวและความร้อน
    \item เพิ่มความสามารถในการอุ้มน้ำและดูดซับธาตุอาหารประจุบวก (CEC สูง)
    \item ลดการชะล้างปุ๋ยไนเตรตและลดการปล่อยก๊าซไนตรัสออกไซด์ได้ถึง 30--50\%
\end{itemize}

\begin{agribox}{การพัฒนาหัวเชื้อจุลินทรีย์ทนแล้งด้วยสารออสโมโพรเทกแทนท์}
ในพื้นที่เกษตรที่ประสบภัยแล้งบ่อยครั้ง การคัดเลือกแบคทีเรีย PGPR ที่สามารถสังเคราะห์สารควบคุมแรงดันออสโมติกภายในเซลล์ (Osmoprotectants) เช่น \textbf{ทรีฮาโลส (Trehalose)} และกรดอะมิโนโพรลีน (Proline) ตลอดจนแบคทีเรียที่ผลิตเมือกเอกโซพอลิแซ็กคาไรด์ (Exopolysaccharides: EPS) ห่อหุ้มรากพืช จะช่วยรักษาความชื้นรอบรากและส่งสัญญาณให้พืชปิดปากใบอย่างสมดุล ส่งผลให้พืชยืนต้นรอดพ้นช่วงวิกฤตแล้งได้ดีกว่าปกติ
\end{agribox}

\section*{คำถามทบทวนท้ายบทที่ 8}
\addcontentsline{toc}{section}{คำถามทบทวนท้ายบทที่ 8}
\begin{enumerate}
    \item ซากเซลล์จุลินทรีย์ (Microbial Necromass) มีบทบาทอย่างไรในการกักเก็บคาร์บอนระยะยาวในดิน
    \item การจัดการน้ำในแปลงนาข้าวแบบเปียกสลับแห้ง (Alternate Wetting and Drying: AWD) ช่วยลดการปล่อยก๊าซมีเทนได้อย่างไรในเชิงจุลชีววิทยา
    \item โครงสร้างรูพรุนของถ่านชีวภาพ (Biochar) มีประโยชน์ต่อการอยู่รอดของแบคทีเรีย PGPR ในดินทรายอย่างไร
\end{enumerate}

\clearpage